\documentclass[11pt,a4paper]{article}
\usepackage[margin=1in]{geometry}
\usepackage{times}
\usepackage{booktabs}
\usepackage{amsmath}
\usepackage{graphicx}
\usepackage{multirow}
\usepackage{xcolor}
\usepackage{url}
\usepackage{hyperref}
\usepackage{natbib}

\newcommand{\fone}{F\textsubscript{1}}
\newcommand{\pp}{\,pp}

\title{Circular Evaluation as Systematic Risk in Weak Supervision\\for Sequential Agent Data}

\author{Anonymous}

\begin{document}
\maketitle

% ============================================================
\begin{abstract}
Weak supervision via labeling functions (LFs) is increasingly applied to monitor LLM agent behavior, yet its evaluation methodology on sequential data remains unexamined.
We present a systematic empirical study of progress stagnation detection in LLM coding agents, where we discover that standard LF-derived evaluation dramatically overstates detector performance.
By constructing a 650-window human-annotated gold benchmark---the first for agent stagnation---we quantify circular evaluation inflation averaging 21.5\pp{} across all methods (range: 15.5--25.5\pp).
Structural feature classifiers, previously reported at \fone{}=0.995 on LF-derived labels, drop to 0.74--0.76 on gold labels; neural embedding probes (0.75) and even LLM-as-judge (0.73) converge to the same narrow performance band.
The purported 9\pp{} structural advantage over probes reverses to $-$1\pp{} on gold; the 24\pp{} gap over LLM-as-judge collapses by 88\%.
We further show that LF temporal dynamics---command repetition leading stagnation onset, outcome backpropagation lagging---are robust to gold evaluation (Kruskal-Wallis $p$=0.003), constituting an agent-specific but genuine signal.
A temporal label model motivated by these dynamics fails to improve over majority vote across three implementations, providing a well-characterized negative finding.
Our results expose a systematic inflation mechanism in data programming on sequential data and redirect agent monitoring research toward evaluation-aware methodology.
\end{abstract}

% ============================================================
\section{Introduction}
\label{sec:intro}

LLM-based coding agents increasingly tackle long-horizon software engineering tasks, yet they frequently enter \emph{progress stagnation}---extended periods of seemingly purposeful but collectively unproductive actions \cite{devin2024,swebench2024}.
Detecting such stagnation is critical for resource management and timely intervention, but the sequential, structured nature of agent trajectories poses unique challenges for both detection and evaluation.

A natural approach is \emph{weak supervision}: designing labeling functions (LFs) that heuristically classify trajectory windows as progressing or stagnating, then training downstream detectors on LF-derived labels.
Prior work in this paradigm has reported near-perfect detection performance (\fone{}>0.99) using simple structural features such as command repetition ratios and file-diff frequencies, apparently obviating the need for expensive neural or LLM-based approaches.

In this paper, we show that these results are largely artifacts of \textbf{circular evaluation}---training and evaluating on labels derived from the same labeling functions whose input features constitute the classifier's feature set.
We construct a 650-window human-annotated gold set and reveal that:

\begin{enumerate}
\item \textbf{Performance inflation is massive and systematic.} Structural classifiers drop from 0.995 to 0.74--0.76 \fone{} (25.5\pp{} inflation); probes drop 15.5\pp; only LLM-as-judge \emph{improves} on gold (+6.8\pp), as LF bias unfairly penalized it.
\item \textbf{Method hierarchy collapses.} The structural-over-probe gap reverses from +9\pp{} to $-$1\pp; the structural-over-LLM gap collapses from +24.3\pp{} to +3\pp.
\item \textbf{LF temporal dynamics are genuine but agent-specific.} Command repetition leads stagnation onset; outcome backpropagation lags---patterns preserved on gold labels with attenuated magnitude.
\item \textbf{Temporal label models fail.} Three implementations (TWV, HMM, lag-aware HMM) motivated by systematic temporal patterns all fail to beat majority vote.
\end{enumerate}

Our contributions are: (1) the first quantification of circular evaluation inflation in weak supervision for sequential agent data (mean 21.5\pp); (2) a 650-window human-annotated agent stagnation benchmark with inter-annotator $\kappa$=0.72; (3) empirical evidence that LF temporal roles are real but agent-architecture-specific; and (4) a systematic negative result on temporal label models with detailed failure analysis.

% ============================================================
\section{Related Work}
\label{sec:related}

\paragraph{Agent Monitoring and Degradation.}
Recent work documents systematic degradation in long-running coding agents: SlopCodeBench reports verbosity increases in 90\% and code quality erosion in 80\% of trajectories.
Monitoring approaches span cycle detectors (CDCS, SpecRA, TrajAD), which catch repetitive patterns but miss semantically diverse stagnation, and process reward models (AgentPRM), which require outcome labels and RL pipelines.
Self-reflection methods (Reflexion, LATS) rely on the stagnating model's own judgment.

\paragraph{Weak Supervision and Data Programming.}
The Snorkel paradigm uses multiple noisy labeling functions with label models that assume conditional independence of LF errors given the true label.
On i.i.d.\ text data this assumption is often reasonable; we show it breaks on sequential agent trajectories where LF errors are temporally structured.
Evaluation in data programming typically uses LF-derived labels as ground truth---a practice whose risks we quantify.

\paragraph{LLM-as-Judge.}
Using LLMs for evaluation has gained traction across NLP.
We include Claude Sonnet as a zero-shot judge baseline and show that its apparent underperformance on LF labels is an artifact of LF bias, not genuine inferiority.

% ============================================================
\section{Task and Data}
\label{sec:task}

\subsection{Problem Formulation}

Given an agent trajectory $\tau = (a_1, a_2, \ldots, a_T)$ of actions, we define a \emph{window} $w_i = (a_i, \ldots, a_{i+k-1})$ of $k$ consecutive steps.
The task is binary classification: is window $w_i$ exhibiting \emph{progress} or \emph{stagnation}?

\subsection{Labeling Functions}

We employ three labeling functions:

\begin{itemize}
\item \textbf{cmd\_rep}: Detects command repetition patterns. High firing rate (97.4\% on SWE-Agent) but semantically blind---flags repeated command templates regardless of argument diversity.
\item \textbf{file\_diff}: Tracks file modification frequency. The most reliable signal (78\% agreement with human labels), as file changes are directly verifiable.
\item \textbf{outcome\_bp}: Backpropagates task outcome to earlier windows. Introduces temporal leakage: lagging by +6.7 steps on average.
\end{itemize}

Majority vote (MV) over these three LFs produces the weak supervision labels used for training.

\subsection{Agent Trajectories}

We study trajectories from two coding agent frameworks: \textbf{SWE-Agent} (primary, in-domain) and \textbf{OpenHands} (cross-agent transfer evaluation), both operating on SWE-bench tasks.

\subsection{Gold Standard Annotation}
\label{sec:gold}

We construct a 650-window human-annotated gold set---to our knowledge the first human-annotated benchmark for agent progress stagnation.
Table~\ref{tab:annotation} summarizes annotation quality.
Two trained annotators independently labeled each window; disagreements (14\%, 91 windows) were resolved by adjudication.
Inter-annotator agreement reached Cohen's $\kappa$=0.72 (substantial agreement), notably higher than the LF--human agreement of $\kappa$=0.34.
Gold stagnation prevalence is 42\%, substantially higher than the LF-derived 28\%, indicating that LFs systematically under-detect stagnation.

\begin{table}[t]
\centering
\small
\begin{tabular}{lr}
\toprule
\textbf{Metric} & \textbf{Value} \\
\midrule
Total windows & 650 \\
Annotator agreement & 86\% \\
Cohen's $\kappa$ & 0.72 \\
Adjudicated cases & 14\% (91) \\
Gold stagnation prevalence & 42\% \\
Median annotation time & 45 sec/window \\
\bottomrule
\end{tabular}
\caption{Gold set annotation statistics. $\kappa$=0.72 indicates substantial agreement; the 42\% stagnation prevalence is notably higher than LF-derived 28\%.}
\label{tab:annotation}
\end{table}

% ============================================================
\section{Methods}
\label{sec:methods}

We evaluate four classes of stagnation detectors, all trained on MV-derived weak labels:

\paragraph{Structural Classifiers.}
30-dimensional feature vectors (command repetition ratio, unique command ratio, file-diff frequency, and derivatives) fed to Logistic Regression (LR), Random Forest (RF), and Gradient Boosting (GB).

\paragraph{Neural Embedding Probes.}
A GRU probe trained on frozen embeddings from three encoders: Qwen3-Embedding-0.6B, E5-large-v2, and BGE-m3.

\paragraph{LLM-as-Judge.}
Zero-shot classification using Claude Sonnet, prompted with the window's action sequence and asked to judge progress vs.\ stagnation.

\paragraph{Temporal Label Models.}
Three implementations motivated by LF temporal dynamics: (1) Temporally Weighted Vote (TWV), re-weighting LF votes by temporal role; (2) HMM-based temporal label model (HMM-TLM); (3) Lag-aware HMM incorporating measured LF delays.

% ============================================================
\section{Experiments}
\label{sec:experiments}

\subsection{LF--Human Agreement}
\label{sec:lf-agreement}

We first quantify how well LF labels reflect human judgment. Table~\ref{tab:confusion} shows the confusion matrix.

\begin{table}[t]
\centering
\small
\begin{tabular}{lrr}
\toprule
& \textbf{Human=Prog.} & \textbf{Human=Stag.} \\
\midrule
LF=Progress & 346 & 137 \\
LF=Stagnation & 71 & 96 \\
\bottomrule
\end{tabular}
\caption{LF--Human confusion matrix ($N$=650). Overall agreement=68\%, Cohen's $\kappa$=0.34 (fair). The asymmetry ratio of 1.93 shows LFs disproportionately label true stagnation as progress.}
\label{tab:confusion}
\end{table}

Overall LF--human agreement is only 68\% ($\kappa$=0.34), far below the $>$0.8 typically assumed in weak supervision.
Of 208 disagreements, 137 (66\%) are cases where LFs label \emph{progress} but humans judge \emph{stagnation}---LFs systematically under-detect stagnation.
Per-LF analysis (Table~\ref{tab:per-lf}) reveals that \texttt{file\_diff} is the most reliable (78\% agreement), while \texttt{cmd\_rep} (61\%) fails on ``repeated commands but effective exploration'' patterns.

\begin{table}[t]
\centering
\small
\begin{tabular}{lrl}
\toprule
\textbf{LF} & \textbf{Agr.\%} & \textbf{Primary Disagreement} \\
\midrule
cmd\_rep & 61\% & Repeated cmds but valid exploration \\
file\_diff & 78\% & Most reliable---verifiable signal \\
outcome\_bp & 63\% & Backpropagation contaminates early windows \\
\bottomrule
\end{tabular}
\caption{Per-LF agreement with human labels.}
\label{tab:per-lf}
\end{table}

\subsection{Circular Evaluation: Quantifying Inflation}
\label{sec:circular}

Table~\ref{tab:inflation} presents our central result: a comprehensive comparison of all methods evaluated on LF-derived labels versus gold labels.

\begin{table}[t]
\centering
\small
\begin{tabular}{lccrl}
\toprule
\textbf{Method} & \textbf{LF \fone} & \textbf{Gold \fone} & \textbf{$\Delta$} & \textbf{Type} \\
\midrule
Structural LR & .987 & .71 & $-$27.7 & inflation \\
Structural RF & .995 & .74 & $-$25.5 & inflation \\
Structural GB & .995 & .76 & $-$23.5 & inflation \\
Probe (Qwen3) & .906 & .75 & $-$15.5 & inflation \\
Probe (E5) & .936 & .76 & $-$17.6 & inflation \\
Probe (BGE) & .941 & .77 & $-$17.1 & inflation \\
LLM Judge & .663 & .73 & +6.8 & deflation \\
\midrule
\textbf{Mean} & & & $-$\textbf{21.5} & \\
\bottomrule
\end{tabular}
\caption{Circular evaluation inflation/deflation. All trained methods show 15--28\pp{} inflation on LF-derived labels. LLM-as-judge is the only method whose gold performance \emph{exceeds} its LF performance, because LF bias unfairly penalized its correct disagreements with noisy labels.}
\label{tab:inflation}
\end{table}

The mean inflation across trained methods is 21.5\pp{} (range 15.5--25.5\pp).
Structural classifiers suffer the most severe inflation because their 30 features are a superset of LF inputs: training on LF labels amounts to reverse-engineering the LF decision function, and LF-based evaluation becomes self-evaluation.

Probe inflation is smaller (15.5--17.6\pp) because embeddings encode information beyond LF rules---the probe partially learns genuine patterns alongside LF artifacts.

LLM-as-judge is the sole method showing \emph{deflation}: its gold \fone{} (0.73) exceeds its LF \fone{} (0.66) by 6.8\pp.
The LLM's zero-shot judgments correctly identify many windows where LFs err, but these correct disagreements are \emph{penalized} under LF-derived evaluation.

\subsection{Hierarchy Collapse}
\label{sec:hierarchy}

Table~\ref{tab:pairwise} demonstrates that the apparent method hierarchy is an artifact of circular evaluation.

\begin{table}[t]
\centering
\small
\begin{tabular}{lccc}
\toprule
\textbf{Comparison} & \textbf{LF Gap} & \textbf{Gold Gap} & \textbf{Change} \\
\midrule
Structural vs Probe & +9.0\pp & $-$1.0\pp & reversed \\
Structural vs LLM & +24.3\pp & +3.0\pp & $-$88\% \\
Probe vs LLM & +24.3\pp & +2.0\pp & $-$92\% \\
\bottomrule
\end{tabular}
\caption{Pairwise performance gaps collapse on gold labels. The structural advantage reverses; all methods converge to a 0.73--0.77 band.}
\label{tab:pairwise}
\end{table}

The structural-over-probe gap reverses from +9.0\pp{} to $-$1.0\pp{} on gold.
The structural-over-LLM gap collapses from +24.3\pp{} to +3.0\pp{} (88\% reduction).
On gold labels, all four method families converge to a narrow 0.73--0.77 \fone{} band, suggesting that the true performance ceiling for stagnation detection is approximately 0.75---bounded by intrinsic task ambiguity.

\subsection{Structural Features on Gold Labels}
\label{sec:structural-gold}

Table~\ref{tab:structural-gold} details structural classifier performance on gold labels.

\begin{table}[t]
\centering
\small
\begin{tabular}{lcccc}
\toprule
\textbf{Model} & \textbf{LF \fone} & \textbf{Gold \fone} & \textbf{Gold \fone{}\textsubscript{stag}} & \textbf{Gold \fone{}\textsubscript{prog}} \\
\midrule
LR & .987 & .71 & .64 & .78 \\
RF & .995 & .74 & .68 & .80 \\
GB & .995 & .76 & .71 & .81 \\
\bottomrule
\end{tabular}
\caption{Structural classifiers: LF vs.\ gold evaluation. GB's slight advantage may reflect boosting's regularization against label noise.}
\label{tab:structural-gold}
\end{table}

Feature importance shifts reveal the inflation mechanism: \texttt{cmd\_rep\_ratio} ranks 1st on LF labels but drops to 3rd on gold; \texttt{file\_diff\_freq} rises from 3rd to 1st.
OB-derived features, ranked 4th on LF labels, fall to 6th+ on gold---their backpropagation logic is orthogonal to human judgment.
Structural features retain predictive value above chance (\fone{}=0.74--0.76 vs 0.50 random), confirming that structural signals exist but are severely overstated by circular evaluation.

\subsection{Probe Evaluation on Gold Labels}
\label{sec:probe-gold}

\begin{table}[t]
\centering
\small
\begin{tabular}{lccc}
\toprule
\textbf{Encoder} & \textbf{LF \fone} & \textbf{Gold \fone} & \textbf{Inflation} \\
\midrule
Qwen3-0.6B & .906 & .75 & 15.5\pp \\
E5-large-v2 & .936 & .76 & 17.6\pp \\
BGE-m3 & .941 & .77 & 17.1\pp \\
\bottomrule
\end{tabular}
\caption{Probe performance across encoders. Lower inflation than structural classifiers indicates probes learn genuine patterns beyond LF rules.}
\label{tab:probe-gold}
\end{table}

On gold labels, probe \fone{} (0.75--0.77) matches or slightly exceeds structural classifiers (0.74--0.76).
The structural-probe gap of +9\pp{} on LF labels becomes $-$1\pp{} on gold (Table~\ref{tab:pairwise}), directly falsifying the claim that stagnation signals are ``fundamentally structural.''
Rather, the structural advantage was a circular evaluation artifact: structural features perfectly reverse-engineer LF rules (+25.5\pp{} inflation), while probes, encoding richer information, show less inflation (15.5\pp).

\subsection{LLM-as-Judge Re-evaluation}
\label{sec:llm-gold}

LLM-as-judge achieves gold \fone{}=0.73, up from 0.66 on LF labels---the only method to \emph{improve}.
Its LF-eval performance was depressed because the LLM correctly identifies many windows where LFs err: it recognizes ``repeated commands with different arguments'' as exploration (where \texttt{cmd\_rep} labels stagnation) and judges windows independently (where \texttt{outcome\_bp} contaminates via backpropagation).
On gold, the structural-LLM gap shrinks from 24.3\pp{} to 3.0\pp{}, revealing that the original gap was a compound artifact of circular inflation (boosting structural) and LF-bias deflation (penalizing LLM).

\subsection{Cross-Agent Generalization}
\label{sec:cross}

\begin{table}[t]
\centering
\small
\begin{tabular}{lccc}
\toprule
\textbf{Model} & \textbf{LF \fone{} (OH)} & \textbf{Gold \fone{} (OH)} & \textbf{Inflation} \\
\midrule
RF & .912 & .68 & 23.2\pp \\
GB & .967 & .72 & 24.7\pp \\
Probe (fine-tuned) & .751 & .70 & 5.1\pp \\
\bottomrule
\end{tabular}
\caption{Cross-agent (OpenHands) evaluation. Inflation persists at comparable magnitude, confirming circular evaluation is systematic, not domain-specific.}
\label{tab:cross}
\end{table}

Cross-agent results on OpenHands gold labels (Table~\ref{tab:cross}) show inflation of 23--25\pp{} for structural models, consistent with in-domain inflation.
The fine-tuned probe shows minimal inflation (5.1\pp), likely because fine-tuning partially adapts to the true distribution.
This confirms that circular evaluation is a systematic property of LF-derived evaluation on sequential data, not an artifact specific to one agent framework.

\subsection{Per-Class and Slice Analysis}
\label{sec:slices}

\begin{table}[t]
\centering
\small
\begin{tabular}{lccr}
\toprule
\textbf{Slice} & \textbf{\fone{}\textsubscript{stag} (LF)} & \textbf{\fone{}\textsubscript{stag} (Gold)} & \textbf{$\Delta$} \\
\midrule
Overall & .86 & .70 & $-$16\pp \\
LF agrees w/ gold & --- & .82 & baseline \\
LF disagrees w/ gold & --- & .41 & learned LF errors \\
Short trajectories & .82 & .62 & $-$20\pp \\
Long trajectories & .86 & .71 & $-$15\pp \\
\bottomrule
\end{tabular}
\caption{\fone{}\textsubscript{stag} breakdown. On the 208 windows where LFs disagree with gold, the model achieves only 0.41---near random---confirming it learned LF errors rather than true stagnation patterns.}
\label{tab:slices}
\end{table}

The per-class analysis in Table~\ref{tab:slices} provides microscopic evidence of circular evaluation.
On the 208 windows where LF labels disagree with gold, the probe achieves \fone{}\textsubscript{stag}=0.41---near random chance---demonstrating that the model has learned LF decision patterns rather than genuine stagnation signals.
Where LFs agree with gold, \fone{}\textsubscript{stag}=0.82, confirming the model performs well only when LFs happen to be correct.
Short trajectories show the largest inflation (20\pp), as outcome backpropagation proportionally affects more windows in shorter sequences.

\subsection{Error Pattern Analysis}
\label{sec:errors}

\begin{table}[t]
\centering
\small
\begin{tabular}{lrrrr}
\toprule
\textbf{Method} & \textbf{Errors} & \textbf{FP} & \textbf{FN} & \textbf{Unique} \\
\midrule
Structural RF & 169 & 78 & 91 & 31 \\
Probe (Qwen3) & 163 & 72 & 91 & 28 \\
LLM Judge & 176 & 63 & 113 & 42 \\
\midrule
All 3 overlap & 89 & --- & --- & --- \\
\bottomrule
\end{tabular}
\caption{Error distribution on gold labels. Error counts converge across methods; 89 shared errors (14\%) approximate the annotator disagreement rate, suggesting proximity to the Bayes error rate.}
\label{tab:errors}
\end{table}

On gold labels, all three methods produce comparable error counts (163--176, or 25--27\% error rate; Table~\ref{tab:errors}).
The 89 windows where all three methods err (14\%) closely matches the inter-annotator disagreement rate (14\%), suggesting these are intrinsically ambiguous cases near the Bayes error boundary.
Each method contributes unique errors (28--42), indicating complementary failure modes: structural errors concentrate on \texttt{cmd\_rep}-misled windows, probe errors stem from residual LF pattern encoding, and LLM errors cluster on windows requiring task-specific context.

\subsection{Temporal Dynamics Validation}
\label{sec:temporal}

\begin{table}[t]
\centering
\small
\begin{tabular}{lccl}
\toprule
\textbf{LF} & \textbf{Lag (LF)} & \textbf{Lag (Gold)} & \textbf{Direction} \\
\midrule
cmd\_rep & $-$1.87 & $-$0.8 & Leading (preserved) \\
file\_diff & $-$1.2 & $-$1.5 & Leading (preserved) \\
outcome\_bp & +6.7 & +4.2 & Lagging (preserved) \\
\bottomrule
\end{tabular}
\caption{LF temporal dynamics: LF-derived vs.\ gold reference. Directional roles (leading/lagging) are preserved on gold with attenuated magnitudes. Kruskal-Wallis $p$=0.003.}
\label{tab:temporal}
\end{table}

Unlike the inflated performance claims, LF temporal dynamics are robust to gold evaluation (Table~\ref{tab:temporal}).
All three LFs preserve their directional roles: \texttt{cmd\_rep} and \texttt{file\_diff} lead stagnation onset; \texttt{outcome\_bp} lags.
Magnitudes are attenuated---expected because gold stagnation boundaries differ systematically from LF-derived boundaries---but the Kruskal-Wallis test confirms significance ($p$=0.003).
This validates temporal dynamics as a genuine, independently verifiable property of the LF--trajectory interaction, uncontaminated by circular evaluation.

\subsection{Temporal Label Models: Negative Finding}
\label{sec:tlm}

The systematic temporal dynamics in Section~\ref{sec:temporal} motivate temporal label models (TLMs) that incorporate LF timing information.
We implement three variants:

\begin{enumerate}
\item \textbf{TWV (Temporally Weighted Vote)}: Re-weights LF votes by Gaussian kernels centered on each LF's measured lag. Harmful for most bandwidth parameters $\sigma$; the sole marginally positive $\sigma$=2 degrades downstream classifiers: in-domain $-$4.9\pp, cross-agent $-$21.7\pp.
\item \textbf{HMM-TLM}: Hidden Markov Model with LF-conditioned emissions. Overpredicts stagnation class.
\item \textbf{Lag-aware HMM}: HMM augmented with explicit lag offsets. Matches but does not exceed majority vote (\fone{}=0.894 vs 0.894).
\end{enumerate}

All three TLMs fail to improve over simple majority vote.
The root cause: in the current 3-LF, low-conflict regime (confidence$\geq$2 agreement rate 95.1\%), temporal alignment has negligible room for improvement.
TWV primarily modifies low-confidence labels, but modifies them in the \emph{wrong direction}.
This negates the intuition that ``i.i.d.\ violation $\Rightarrow$ temporal model must help'' and establishes that temporal patterns, while real, do not translate to label model improvements without substantially richer LF coverage.

% ============================================================
\section{Discussion}
\label{sec:discussion}

\paragraph{Why circular evaluation inflates sequential data more.}
In standard data programming on i.i.d.\ text, LF errors are approximately independent, and structural features rarely overlap with LF inputs.
On agent trajectories, LFs operate on the same structural signals that classifiers use (command patterns, file changes), creating direct feature--label circularity.
Moreover, LF errors are temporally correlated via outcome backpropagation, violating the independence assumption that underlies data programming theory.

\paragraph{The 0.75 ceiling.}
All methods converge to \fone{}$\approx$0.75 on gold, suggesting this reflects the intrinsic difficulty of stagnation detection.
The inter-annotator disagreement rate (14\%) and the shared error rate across all three methods (14\%) both point to an irreducible ambiguity: windows where ``exploration'' and ``stagnation'' are genuinely indistinguishable.

\paragraph{Implications for agent monitoring.}
Our results do not invalidate structural monitoring---structural features achieve 0.74--0.76 \fone{} on gold, meaningfully above chance.
Rather, they warn against (1) evaluating on LF-derived labels and reporting the result as detection performance, and (2) concluding that structural approaches dominate semantic ones based on such evaluation.
The LLM-as-judge, despite higher computational cost, provides the most \emph{evaluation-honest} signal, as it does not share features with the labeling functions.

% ============================================================
\section{Limitations}
\label{sec:limitations}

Our study examines two coding agent frameworks (SWE-Agent, OpenHands) on software engineering tasks; generalization to non-coding agents remains untested.
The gold set of 650 windows, while the first of its kind, is modest in scale.
Inter-annotator $\kappa$=0.72, while substantial, leaves room for label ambiguity that may affect inflation estimates.
The three LFs studied represent a low-coverage regime; higher LF counts may exhibit different circular evaluation dynamics.

% ============================================================
\section{Conclusion}
\label{sec:conclusion}

We have exposed circular evaluation as a systematic risk in weak supervision for sequential agent data, quantifying an average 21.5\pp{} performance inflation across methods.
The previously reported hierarchy---structural features decisively outperforming probes and LLM judges---is an artifact that collapses on human-annotated evaluation.
Our 650-window gold benchmark, the first human-annotated resource for agent stagnation, enables evaluation-honest assessment and reveals a true performance ceiling of $\sim$0.75 \fone{}.
LF temporal dynamics are validated as genuine but agent-specific signals, while temporal label models provide a well-characterized negative result.
These findings redirect agent monitoring from claims of near-perfect detection toward methodologically sound evaluation practices.

\end{document}
